% !TEX root = /Users/nai1ka/Projects/thesis_template/thesis.tex
\chapter{Conclusion}
\label{chap:conclusion}

This thesis set out to develop a system for continuous performance monitoring and automated regression detection for Android applications. This chapter summarises the work, answers the research questions, states the main contribution, and outlines directions for future work.

\section{Summary of the Work}
\label{sec:concl_summary}

This thesis has addressed the problem of detecting performance regressions in Android applications after release. Mobile applications run on a wide range of devices, and a change that keeps the application functionally correct can still make it slower or less responsive for users. Pre-release profiling tools cannot reproduce this variety, while existing post-release monitoring tools either provide no automated alerting or rely on fixed thresholds that must be tuned by hand. The thesis identified the gap between continuous post-release monitoring and automated regression detection, and set out to build a system that closes it.

To address this gap, a performance monitoring system was designed and implemented. The system consists of four components. A mobile SDK, written in Kotlin, collects performance metrics on real user devices and sends them to a backend in batches. The backend, built with Ktor, receives the metrics, assigns each device to a performance cohort, and stores the data in ClickHouse for time-series metrics and in PostgreSQL for user and project information. A regression detection pipeline compares a recent current window against a historical baseline window, using a P95 percentile shift to find candidate regressions and a Mann-Whitney U test to confirm them. A frontend dashboard, built with Streamlit, lets developers manage projects, explore metrics, and configure alerts.

The system was validated in three steps, as described in Chapter~\ref{chap:eval}. The overhead of the SDK was measured on real devices, the accuracy of the collected metrics was checked against reference tools, and the regression detection pipeline was evaluated on synthetic regressions.

\section{Answers to the Research Questions}
\label{sec:concl_rq}

The thesis set out to answer three research questions.

The first research question asked how a low-overhead, continuous performance monitoring system for Android applications can be implemented so that it runs on real user devices without significantly affecting the user experience. This question was answered by the design and implementation of the mobile SDK (Chapters~\ref{chap:met} and~\ref{chap:impl}). The SDK collects CPU usage, memory usage, frame time, startup time, interaction delay, and ANR events through dedicated collectors, stores the data in a two-tier buffer, and transmits it to the backend in batches to reduce the number of network connections and disk writes. The overhead evaluation in Section~\ref{sec:eval_overhead} measured an average CPU overhead of [TODO: value]\% and a memory overhead of [TODO: value]~MB. [TODO: state whether these results confirm that the SDK is low-overhead enough to run on real devices without significantly affecting the user experience.]

The second research question asked how performance regressions can be detected accurately using statistical methods while keeping the number of false alerts low. This question was answered by the regression detection pipeline (Section~\ref{sec:detection} and Section~\ref{sec:regression_impl}). The pipeline partitions the data by metric, application screen, and device cohort. For each partition, the main filter is a rolling-window P95 percentile shift, which reports a candidate regression only when the 95th percentile of a metric increases by more than the configured threshold. A Mann-Whitney U test is then applied as a secondary check that removes changes likely to be caused by random noise. In the control run reported in Section~\ref{sec:eval_regression}, the pipeline produced [TODO: number] confirmed false alerts. [TODO: comment on this number as evidence that the method helps to reduce false alerts; do not claim that false alerts are eliminated.]

The third research question asked how effective the proposed system is at detecting performance regressions when tested on real Android applications. This question was answered by the regression detection evaluation in Section~\ref{sec:eval_regression}. The pipeline was tested with three types of synthetic regression — increased CPU load, a memory leak, and UI thread blocking — and detected [TODO: number] of the [TODO: number] injected regressions. [TODO: add one sentence interpreting this result.] Because the regressions were synthetic, this result shows that the detection method works as designed under controlled conditions; its effectiveness against real regressions in production has not yet been measured.

\section{Main Contribution}
\label{sec:concl_contribution}

The main contribution of this thesis is a performance monitoring system that combines continuous, low-overhead post-release monitoring of Android applications with automated, server-side performance regression detection. Earlier work usually addresses only one of these two parts. Pre-release profilers provide detailed analysis under controlled conditions, post-release monitoring tools collect data from real devices but commonly provide no automated alerting or rely on fixed thresholds, and statistical regression detection methods are mostly studied for large-scale server systems rather than for mobile applications. The contribution of this work is the integration of these parts into a single system that:

\begin{itemize}
      \item collects performance metrics from real user devices with low overhead;
      \item groups the collected data by metric, application screen, and device performance cohort, so that devices with different hardware are not compared directly;
      \item stores large volumes of time-series telemetry efficiently;
      \item compares recent performance against a historical baseline using a P95 percentile shift;
      \item confirms candidate regressions with a non-parametric statistical test before alerting;
      \item notifies developers automatically when a regression is confirmed.
\end{itemize}

This combination removes the need for developers to set and tune fixed thresholds by hand, and it reports regressions that are confirmed statistically rather than triggered by a single measurement. The system is not intended to replace pre-release profiling tools, which still provide more detailed analysis during development. Instead, it complements them by covering the period after release, when the application runs on real devices under real usage conditions.

\section{Future Work}
\label{sec:concl_future}

Several directions remain for future work.

The most important next step is to evaluate the system in a real deployment. The current evaluation used synthetic regressions on a small set of test devices. Deploying the SDK in an application with real users would show how the system behaves under varied hardware, usage patterns, and network conditions, and would allow its effectiveness to be measured against real regressions rather than injected ones.

The set of collected metrics could also be extended. The current SDK collects CPU usage, memory usage, and several UI metrics. Battery consumption and network usage are also important for mobile applications~\cite{Saraiva}, and adding collectors for them would give a more complete picture of application performance.

The regression detection method could be made more adaptive. At present, the relative degradation threshold is configured per metric and per screen. Future work could replace this fixed threshold with a method that derives a suitable threshold from the historical data of each cohort, which would further reduce the need for manual configuration. The detection pipeline could also be extended to account for recurring daily or weekly usage patterns, which the current fixed-window comparison does not capture.

The device cohort classification could be refined. Cohorts are currently assigned from total RAM and CPU core count. Additional hardware and software characteristics, such as the CPU model or the Android version, could produce cohorts that group devices with similar real-world performance more accurately.

Finally, the system could be extended to help developers find the cause of a regression. The backend already stores the application version with each metric. Future work could link a confirmed regression to the application version in which it first appeared, which would help developers connect a regression to a specific release.
